% related-work.tex -- EVM Executor Paper (BNR-2026-008)

\section{Related Work}

\subsection{zkEVM Type Spectrum}

Buterin~\cite{buterin2022zkevm} introduced the zkEVM type classification
ranging from Type~1 (fully Ethereum-equivalent) to Type~4 (language-compatible
only). Chaliasos et al.~\cite{chaliasos2024analyzing} formalized this spectrum
and benchmarked production zkEVM systems, finding that Type~2--3 systems
(Polygon zkEVM, Scroll) achieve practical proving times of 190--200 seconds
per batch while maintaining high EVM compatibility. The constraint-level
analysis by Burchardt et al.~\cite{burchardt2025constraint} demonstrates that
PLONKish arithmetizations with custom gates reduce constraint counts by
$5\times$--$30\times$ compared to R1CS for arithmetic-heavy opcodes.

\subsection{Production zkEVM Architectures}

Five production systems inform the architectural decisions in this paper.

\textbf{Polygon CDK (cdk-erigon).} Polygon migrated from a custom Go
implementation (zkevm-node, archived February 2025) to an Erigon
fork~\cite{polygon2024cdk}. The execution layer uses Go (Erigon), while the
proving layer is a separate C++ executor that processes transactions through
zkASM (a custom assembly language). This separation achieves 10$\times$ less
disk usage and 150$\times$ faster synchronization compared to the original
zkevm-node. The key insight is that \emph{the proving executor is architecturally
separate from the EVM executor}.

\textbf{Scroll (scroll-geth).} Scroll maintains a direct fork of
go-ethereum~\cite{scroll2024whitepaper} that replaces the Merkle Patricia Trie
(MPT) with a ZK-friendly Poseidon-based trie (zktrie), activated via
\texttt{config.scroll.useZktrie}. However, the 2025 Euclid
upgrade~\cite{scroll2025euclid} migrates \emph{back} to standard MPT combined
with OpenVM (a general-purpose zkVM), suggesting that tightly coupling the
state trie to the EVM fork creates unsustainable maintenance burden.

\textbf{zkSync Era (EraVM).} Rather than forking Geth, zkSync built a custom
register-based virtual machine (EraVM)~\cite{zksync2024architecture} optimized
for ZK circuits. An EVM bytecode interpreter was added on top for compatibility.
This approach trades compatibility (Type~4) for proving efficiency, with the
Boojum~2.0 upgrade targeting native EVM execution.

\textbf{Optimism (op-geth).} Optimism's Geth fork~\cite{optimism2024opgeth}
demonstrates the minimal viable diff: 16{,}881 lines added and 664 deleted
against upstream Geth, with only 34~lines touching the EVM itself (precompile
overrides, caller overrides). While op-geth is an optimistic rollup (no ZK
proving), it establishes an upper bound on fork minimality.

\textbf{Reth.} Paradigm's Rust EVM implementation~\cite{paradigm2024reth}
achieves 100--200~MGas/s in live sync and up to 1.5~GGas/s with parallel
execution (Gravity Reth)~\cite{galxe2025gravity}. The revmc JIT
compiler~\cite{paradigm2024revmc} further improves interpreted EVM performance
by 1.85$\times$--19$\times$. These benchmarks establish the performance ceiling
for EVM execution engines.

\subsection{EVM Tracing Infrastructure}

Geth provides three tracing approaches with different performance
characteristics. The \texttt{structLogger} captures every opcode with full
stack, memory, and storage state, incurring 10--50$\times$
overhead~\cite{geth2024tracing}. The \texttt{callTracer} operates at call-frame
granularity with lower overhead. Since Geth~1.14+, the
\texttt{core/tracing} hooks API~\cite{geth2024hooks} provides event-driven
callbacks (\texttt{OnOpcode}, \texttt{OnStorageChange},
\texttt{OnBalanceChange}, \texttt{OnNonceChange}) that enable selective
capture of only the operations relevant to ZK witness generation. Production
zkEVMs universally use custom compiled tracers rather than the built-in
JavaScript-based tracing infrastructure.

\subsection{ZK Constraint Costs}

The Polygon zkEVM ROM specification~\cite{polygon2024zkcounters} documents
per-opcode ZK counter costs across five resource dimensions: arithmetic
(\texttt{cnt\_arith}), binary (\texttt{cnt\_binary}), Keccak-f permutations
(\texttt{cnt\_keccak\_f}), Poseidon operations (\texttt{cnt\_poseidon\_g}),
and memory alignment (\texttt{cnt\_mem\_align}). KECCAK256 dominates with
192~arithmetic counters, 193~binary counters, and 2~Keccak-f permutations,
translating to approximately 150{,}000 R1CS constraints per
invocation~\cite{burchardt2025constraint}. Storage operations (SLOAD, SSTORE)
require 255~Poseidon operations each for state trie traversal.
